\documentclass[11pt]{article}
\usepackage[margin=1in]{geometry}
\usepackage{amsmath,amssymb,amsthm,mathtools,mathrsfs,bm}
\usepackage{enumitem}
\usepackage{booktabs}
\usepackage{microtype}
\usepackage{hyperref}

\allowdisplaybreaks

\newtheorem{theorem}{Theorem}[section]
\newtheorem{definition}[theorem]{Definition}
\newtheorem{proposition}[theorem]{Proposition}
\newtheorem{lemma}[theorem]{Lemma}
\newtheorem{corollary}[theorem]{Corollary}
\theoremstyle{remark}
\newtheorem{remark}[theorem]{Remark}

\newcommand{\R}{\mathbb R}
\newcommand{\C}{\mathbb C}
\newcommand{\N}{\mathbb N}
\newcommand{\Z}{\mathbb Z}
\newcommand{\Tset}{\mathbb T}
\newcommand{\mcN}{\mathcal N}
\newcommand{\mcE}{\mathcal E}
\newcommand{\mcL}{\mathcal L}
\newcommand{\one}{\mathbf 1}
\newcommand{\tr}{\operatorname{tr}}

\title{Finite $SU(3)$ Lattice Gauge Theory\\
from Particle Trajectories}
\author{}
\date{}
\begin{document}
\maketitle

\begin{abstract}
We construct a finite, explicit $SU(3)$ lattice gauge theory from the discrete
trajectory data of a finite particle system.
Starting from a recorded trajectory data set, we extract a finite lattice of
sites together with temporal worldline links and same-time interaction links.
At each site we form a complex color triplet by modulating the viscous force
vector with a kinematic phase factor of de~Broglie type; these triplets encode
momentum direction in the internal color space.
Link transport in $SU(3)$ is then defined from short-link color increments,
which makes the transporter asymptotically local in the small-mesh regime.
Canonical plaquette holonomies, discrete field strengths, a finite Wilson action,
and covariant fermionic and scalar matter actions are constructed on the same
lattice.
The complete data package $\mathcal Q$ is proved to be finite and $SU(3)$
gauge-invariant.
In the stationary thermodynamic regime, scaled link and plaquette sums converge
to integrals against deterministic mean-field measures.
Under an admissible mean-field refinement regime with a $C^2$ color-field limit,
shape-regular plaquettes, and plaquette-density convergence, the discrete Wilson
action converges to the continuum Yang--Mills functional
\[
\frac14\int_M \tr(F_{\mu\nu}F^{\mu\nu})\,d\mathrm{vol}_g.
\]
Path-integral quantization and axiom verification are deferred to companion papers.
\end{abstract}

\tableofcontents

\section{Introduction}
\label{sec:intro}

Lattice gauge theory provides
a non-perturbative regularization of quantum field theories by replacing
continuous spacetime with a discrete lattice and assigning $SU(N)$ group elements
to the links.
In the standard formulation the link variables are
\emph{independent integration variables} of a path integral, and the gauge group
is imposed by hand.
The present paper takes a different starting point: we ask whether a natural
$SU(3)$ gauge structure can be \emph{extracted} from discrete trajectory data,
without imposing any gauge group a priori.

The data source is the finite-particle kinetic core studied in the companion
convergence paper~\cite{convergence}.
At each discrete timestep, each particle $i$ experiences a force
\[
  \mathbf F^{\mathrm{visc}}_i
  = \nu_{\mathrm{visc}}\sum_{j\neq i}\omega_{ij}(v_j - v_i)
\]
that pulls its velocity toward a weighted average of its interaction partners.
This force is a vector in $\R^3$ and therefore lives naturally in the fundamental
representation of $O(3)$.

The key observation is that $\R^3$ is three-dimensional, matching the defining
dimension of the fundamental representation of $SU(3)$, the gauge group used
throughout this paper.
By complexifying the viscous force with a kinematic phase factor analogous to the
de~Broglie phase $e^{ip\ell/\hbar}$, each site acquires a complex
color triplet~$\mathbf c_x\in\C^3$.
The color triplets at two neighboring sites define a short-link color increment
from which we build an \(SU(3)\) parallel transport \(U_3(e)\) for every link \(e\).
No additional input is required: the gauge group, the link variables, and
the coupling strength all emerge from the dynamics of the viscous particle system.

In detail, we first build a complex color triplet at each site, then project the
short-link color increment onto the Gell-Mann basis to obtain a Hermitian
traceless edge generator. The coupling parameter \(g_3>0\) is placed to match
the usual continuum normalization in
\[
D_\mu=\partial_\mu-ig_3 A_\mu^a\lambda_a/2.
\]
We then add temporal worldline links and define canonical elementary plaquettes
from one time step and one persistent interaction link. This restores the local
2-cell geometry needed for a genuine Yang--Mills continuum limit without
reintroducing any extra edge nomenclature.

This paper supplies all the formal definitions and proofs required for a
subsequent axiom-verification paper.
The finite-data result is Theorem~\ref{thm:package}: the tuple
\[
\mathcal Q = (\text{lattice},\,\{U_3(e)\},\,S_\mathrm W,\,S_\mathrm F,\,S_\phi)
\]
is finite, explicit, and $SU(3)$ gauge-invariant.
The continuum result is Theorem~\ref{thm:continuum-limit-ym}: under the stated
refinement hypotheses, the same data recover the continuum Yang--Mills action.

\paragraph{Organization.}
Section~\ref{sec:notation} fixes notation, imports kernel parameters from the
companion convergence paper, and tabulates all scalar parameters used throughout.
Section~\ref{sec:run} constructs the finite lattice $\mathcal L=(\mcN,\mcE)$
from the recorded data, including worldline and interaction links.
Section~\ref{sec:forces} defines the viscous force on sites and links and proves
its covariance under orthogonal transformations.
Section~\ref{sec:su3} introduces color triplets, the Gell-Mann basis, and the
$SU(3)$ link transport.
Section~\ref{sec:holonomy} defines path holonomies, canonical plaquettes,
discrete field strengths, the Wilson action, and proves gauge invariance.
Section~\ref{sec:matter} constructs covariant fermionic and scalar matter actions.
Section~\ref{sec:package} packages the complete data object and states the finiteness
theorem.
Sections~\ref{sec:discussion} and~\ref{sec:conclusions} discuss the lattice-gauge
interpretation, the thermodynamic graph and plaquette limits, the geometric
continuum limit to Yang--Mills theory, and quantization.

\section{Notation and setup}
\label{sec:notation}

\begin{definition}[Kernel parameters and trajectory data]
\label{def:imports}
Let $d=3$.
The spatial dimension is fixed to three throughout because the viscous force is
a vector in physical space and its three components form a natural triplet
for the fundamental representation of $SU(3)$.
Fix $N\ge1$ and the index set $\mathcal I:=\{1,\dots,N\}$.
At each time every trajectory entry is recorded for all $i\in\mathcal I$.
The alive set $\mathcal A(t)\subseteq\mathcal I$ records which indices are active.
Fix kernel parameters
\[
\ell_{\mathrm{visc}}>0,\qquad \kappa_0>0,\qquad \nu_{\mathrm{visc}}\ge0,
\]
and define the regularized Gaussian kernel
\[
K_\rho(x,y):=\kappa_0+\exp\!\left(-\frac{\|x-y\|^2}{2\ell_{\mathrm{visc}}^2}\right),
\qquad \rho:=\ell_{\mathrm{visc}}.
\]
The kernel satisfies
\[
\kappa_*:=\inf_{x,y}K_\rho(x,y)=\kappa_0>0,\qquad
\kappa^*:=\sup_{x,y}K_\rho(x,y)=1+\kappa_0<\infty.
\]
These bounds are imported from the companion paper on kinetic convergence;
no other results from that paper are used here.

A recorded trajectory is available at discrete times
\[
\Tset=\{0,1,\dots,T\},\quad T\in\N,
\]
with finite alive sets $\mathcal A(t)\subseteq\mathcal I$ and position,
velocity, and fitness data
\[
x_i(t),\,v_i(t)\in\R^3,\qquad \Phi_i(t)\in\R,\qquad i\in\mathcal I,\ t\in\Tset.
\]
At each timestep $t$, the finite set
$\mathcal P_t\subseteq \mathcal A(t)\times\mathcal A(t)\setminus\{(i,i):i\in\mathcal A(t)\}$
records the ordered interaction companions used by the algorithm.
\end{definition}

The following table summarizes all scalar parameters introduced in this paper,
together with their role and the constraint they must satisfy.

\begin{center}
\begin{tabular}{@{}lll@{}}
\toprule
Symbol & Role & Constraint \\
\midrule
$\ell_{\mathrm{visc}}$ & Viscous kernel length scale & $>0$ \\
$\kappa_0$ & Kernel floor (regularization) & $>0$ \\
$\nu_{\mathrm{visc}}$ & Viscous coupling strength & $\ge0$ \\
$\Delta t$ & Physical time spacing between recorded slices & $>0$ \\
$g_3$ & Gauge coupling constant & $>0$ \\
$\ell_0$ & Length unit for de~Broglie phase & $>0$ \\
$\hbar_{\mathrm{eff}}$ & Effective action quantum & $>0$ \\
$\mu_p$ & Momentum coupling to phase & $>0$ \\
$\epsilon_c$ & Color regularization & $>0$ \\
$\eta$ & Distance regularization in scalar action & $>0$ \\
$m_F$ & Fermion mass & $>0$ \\
$m_\phi$ & Scalar mass & $>0$ \\
\bottomrule
\end{tabular}
\end{center}

\begin{lemma}[Kernel orthogonal symmetry]
\label{lem:kernel-orth}
For every orthogonal matrix $O\in O(3)$ and all $x,y\in\R^3$,
\[
K_\rho(Ox,Oy)=K_\rho(x,y).
\]
\end{lemma}
\begin{proof}
By definition,
\[
\|Ox-Oy\|^2=\|O(x-y)\|^2=\|x-y\|^2,
\]
so the Gaussian factor is invariant under simultaneous orthogonal transformation.
\end{proof}

\begin{lemma}[Orthogonal covariance of viscous forces]
\label{lem:orthogonal}
Let $O\in O(3)$ and define transformed data $x_i'(t)=Ox_i(t)$, $v_i'(t)=Ov_i(t)$.
Then for each site $x=n_{i,t}$,
\[
\mathbf F^{\mathrm{visc}}(x;\text{transformed data})=O\,\mathbf F^{\mathrm{visc}}(x;\text{original data}).
\]
\end{lemma}
\begin{proof}
Lemma~\ref{lem:kernel-orth} gives $K_\rho(Ox_i,Ox_j)=K_\rho(x_i,x_j)$, hence
\[
\deg_{\mathrm{link}}'(i,t)=\deg_{\mathrm{link}}(i,t),
\qquad
\omega_{ij}^{\mathrm{link}\prime}(t)=\omega_{ij}^{\mathrm{link}}(t).
\]
For each link $e=(n_{i,t}\to n_{j,t})$,
\[
\mathbf F^{\mathrm{visc}}{}'(e)
=\nu_{\mathrm{visc}}\,\omega^{\mathrm{link}}_{ij}(t)(Ov_j-Ov_i)
=O\,\mathbf F^{\mathrm{visc}}(e).
\]
Summing outgoing links at a fixed source site preserves left multiplication by $O$.
\end{proof}

\section{Recorded run and lattice}
\label{sec:run}
\begin{definition}[Recorded particle trajectory data]
\label{def:run}
The recorded data set is the tuple
\[
\mathscr R=\Bigl(
N,\Tset,\{\mathcal A(t)\}_{t\in\Tset},\{\mathcal P_t\}_{t\in\Tset},
\{x_i(t),v_i(t),\Phi_i(t)\}_{i\in\mathcal I,t\in\Tset},\{K_\rho\}\Bigr).
\]
\end{definition}

\begin{definition}[Site set and projections]
\label{def:sites}
For each pair $(i,t)\in\mathcal I\times\Tset$ define a lattice site $n_{i,t}$ and set
\[
\mcN:=\{n_{i,t}\mid i\in\mathcal I,\ t\in\Tset\}.
\]
Define physical time, spatial, velocity, and spacetime projections by
\[
\mathbf t(n_{i,t}) := t\,\Delta t,
\]
\[
\mathbf x(n_{i,t})=x_i(t),\qquad \mathbf v(n_{i,t})=v_i(t).
\]
\[
\Xi(n_{i,t}):=\bigl(\mathbf t(n_{i,t}),\mathbf x(n_{i,t})\bigr)\in\R^4.
\]
All maps are well-defined by definition of $\mcN$.
\end{definition}

\begin{definition}[Oriented links]
\label{def:links}
Define the interaction and temporal link sets by
\[
E_{\mathrm{int}}:=\{(n_{i,t}\to n_{j,t}):\,(i,j)\in\mathcal P_t,\ t\in\Tset\},
\]
\[
E_{\mathrm{time}}:=\{(n_{i,t}\to n_{i,t+1}):\,i\in\mathcal I,\ 0\le t<T\}.
\]
For convenience, write
\[
e_{\mathrm{int}}(n_{i,t}\to n_{j,t}) := (n_{i,t}\to n_{j,t}),
\qquad
e_{\mathrm{time}}(n_{i,t}\to n_{i,t+1}) := (n_{i,t}\to n_{i,t+1}).
\]
By construction all source and target sites of $E_{\mathrm{int}}$ and
$E_{\mathrm{time}}$ lie in $\mcN$ because $\mcN$ is defined for all
$(i,t)\in\mathcal I\times\Tset$.
Set
\[
\mcE:=E_{\mathrm{int}}\cup E_{\mathrm{time}}.
\]
The lattice is $\mathcal L:=(\mcN,\mcE)$.
For $e=(u\to v)\in\mcE$, write $s(e):=u$ and $t(e):=v$, define the
spacetime displacement
\[
\xi_e:=\Xi(t(e))-\Xi(s(e))\in\R^4,
\qquad
\ell_e:=\|\xi_e\|,
\]
and define $e^{-1}:=(t(e)\to s(e))$.
\end{definition}

\begin{proposition}[Finiteness and cardinality]
\label{prop:finiteness}
All sets in Definitions~\ref{def:run} and~\ref{def:links} are finite.
With $m_t:=|\mathcal P_t|$,
\[
|E_{\mathrm{int}}|=\sum_{t=0}^{T}m_t,
\qquad
|E_{\mathrm{time}}|=NT,
\qquad
|\mcE|=NT+\sum_{t=0}^{T}m_t.
\]
\end{proposition}
\begin{proof}
Finiteness follows from finiteness of $\Tset$ and all $\mathcal P_t$.
The cardinality identity holds by one-to-one correspondence: each
$(i,j,t)\in\bigcup_{t=0}^{T}\mathcal P_t\times\{t\}$ contributes one interaction
link and each $(i,t)\in\mathcal I\times\{0,\dots,T-1\}$ contributes one temporal
link.
\end{proof}

\section{Viscous force from interaction links}
\label{sec:forces}

\begin{definition}[Link degree and pairwise force coefficients]
\label{def:ig-weights}
For each $(i,t)$ define the raw link degree and normalized weights
\[
\deg_{\mathrm{link}}(i,t):=\sum_{(i,j)\in\mathcal P_t}K_\rho(x_i(t),x_j(t)),
\]
\[
\omega^{\mathrm{link}}_{ij}(t):=
\begin{cases}
\dfrac{K_\rho(x_i(t),x_j(t))}{\deg_{\mathrm{link}}(i,t)},&(i,j)\in\mathcal P_t,\ \deg_{\mathrm{link}}(i,t)>0,\\[0.3em]
0,&\text{otherwise.}
\end{cases}
\]
\end{definition}

\begin{definition}[Viscous pair force and site force]
\label{def:forces}
For each interaction link $e=(n_{i,t}\to n_{j,t})\in E_{\mathrm{int}}$ define
\[
\mathbf F^{\mathrm{visc}}(e):=\nu_{\mathrm{visc}}\omega^{\mathrm{link}}_{ij}(t)\,(v_j(t)-v_i(t))\in\R^3.
\]
For $x=n_{i,t}\in\mcN$ define
\[
\mathbf F^{\mathrm{visc}}(x):=\sum_{\substack{e\in E_{\mathrm{int}}\\ s(e)=x}}\mathbf F^{\mathrm{visc}}(e).
\]
If the sum is over an empty set, it is $0$.
\end{definition}

\begin{lemma}[Weight properties]
\label{lem:weights}
For every $i\in\mathcal A(t)$ and $t\in\Tset$,
\[
\omega^{\mathrm{link}}_{ij}(t)\ge0,\qquad
\sum_{j:(i,j)\in\mathcal P_t}\omega^{\mathrm{link}}_{ij}(t)=\mathbf 1_{\{\deg_{\mathrm{link}}(i,t)>0\}}.
\]
\end{lemma}
\begin{proof}
Positivity is immediate from $K_\rho\ge\kappa_*>0$.
If $\deg_{\mathrm{link}}(i,t)>0$, summing the definitions over $j$ gives
\[
\sum_{j:(i,j)\in\mathcal P_t}\omega^{\mathrm{link}}_{ij}(t)
=\frac{1}{\deg_{\mathrm{link}}(i,t)}\sum_{j:(i,j)\in\mathcal P_t}K_\rho(x_i(t),x_j(t))=1.
\]
If $\deg_{\mathrm{link}}(i,t)=0$, every summand is $0$ by convention.
\end{proof}

\begin{lemma}[Node reconstruction of viscous force]
\label{lem:reconstruction}
For every $x=n_{i,t}\in\mcN$,
\[
\mathbf F^{\mathrm{visc}}(x)=\sum_{\substack{e\in E_{\mathrm{int}}\\ s(e)=x}}\mathbf F^{\mathrm{visc}}(e).
\]
\end{lemma}
\begin{proof}
By Definition~\ref{def:links}, the interaction links with source $n_{i,t}$ in
$E_{\mathrm{int}}$
are exactly $\{(n_{i,t}\to n_{j,t}):(i,j)\in\mathcal P_t\}$, which is the same
index set used in Definition~\ref{def:forces}.
\end{proof}

\begin{proposition}[Orthogonal basis covariance]
\label{thm:viscous-su3}
Let $O\in O(3)$ and consider transformed data
$x_i'(t)=Ox_i(t)$, $v_i'(t)=Ov_i(t)$.
For every site $x\in\mcN$,
\[
\mathbf F^{\mathrm{visc}}(x;\text{transformed data})=O\,\mathbf F^{\mathrm{visc}}(x;\text{original data}).
\]
Moreover, real orthogonal matrices satisfy $O^\top=O^{-1}$, hence $O^\dagger=O^{-1}$,
so $O\in U(3)$.
Consequently, orthogonal changes of the physical coordinate frame embed in
unitary changes of the complexified color index.
\end{proposition}
\begin{proof}
Covariance of the force is Lemma~\ref{lem:orthogonal}.
For the embedding: since $O$ has real entries, $O^\dagger=O^\top=O^{-1}$,
so $O^\dagger O=I$ and $O$ is unitary when viewed as a complex matrix.
\end{proof}

\section{Color triplets and $SU(3)$ link transport}
\label{sec:su3}

\begin{definition}[Normalized color vectors]
\label{def:colors}
Fix constants $g_3>0,\ \ell_0>0,\ \hbar_{\mathrm{eff}}>0,\ \mu_p>0,\ \epsilon_c>0$.
For $x=n_{i,t}\in\mcN$ and each coordinate $\alpha\in\{1,2,3\}$, define
\[
\tilde c_x^{(\alpha)}:=(\mathbf F^{\mathrm{visc}}(x))_\alpha
\exp\!\left(i\frac{\mu_p\,v_i^{(\alpha)}(t)\ell_0}{\hbar_{\mathrm{eff}}}\right),
\]
\[
\mathbf c_x:=\frac{\tilde{\mathbf c}_x}{\sqrt{\tilde{\mathbf c}_x^\dagger\tilde{\mathbf c}_x+\epsilon_c^2}}\in\C^3.
\]
\end{definition}

\begin{theorem}[Emergent local $SU(3)$ color symmetry]
\label{thm:viscous-su3-local}
For any map $\Omega:\mcN\to O(3)$ and each site $x\in\mcN$, define
\[
\mathbf c_x^\Omega:=\Omega(x)\mathbf c_x.
\]
Because $\Omega(x)$ is real orthogonal, $\|\mathbf c_x^\Omega\|=\|\mathbf c_x\|$ for every
$x\in\mcN$.
Hence the color data carry a local $O(3)$ redundancy, which embeds as real
unitaries into $U(3)$.
Imposing $\det\Omega(x)=1$ for every site yields a local $SO(3)$ subgroup
$SO(3)\subset SU(3)$ on each site.
The full local $SU(3)$ bundle symmetry used below is then defined explicitly by
$\Omega:\mcN\to SU(3)$ in Definition~\ref{def:gauge}.
\begin{proof}
\[
\mathbf c_x^\Omega
=\frac{\Omega(x)\tilde{\mathbf c}_x}
{\sqrt{\tilde{\mathbf c}_x^\dagger\tilde{\mathbf c}_x+\epsilon_c^2}}
=\Omega(x)\mathbf c_x.
\]
Since $\Omega(x)$ is orthogonal,
\[
\|\mathbf c_x^\Omega\|^2
=\frac{\tilde{\mathbf c}_x^\dagger\Omega(x)^\dagger\Omega(x)\tilde{\mathbf c}_x}
{\tilde{\mathbf c}_x^\dagger\tilde{\mathbf c}_x+\epsilon_c^2}
=\frac{\tilde{\mathbf c}_x^\dagger\tilde{\mathbf c}_x+\epsilon_c^2}
{\tilde{\mathbf c}_x^\dagger\tilde{\mathbf c}_x+\epsilon_c^2}
=\|\mathbf c_x\|^2.
\]
Therefore the normalization denominator is unchanged and
$\|\mathbf c_x^\Omega\|=\|\mathbf c_x\|$.
\end{proof}
\end{theorem}

\begin{remark}[Physical interpretation of the color construction]
\label{rem:color-physics}
The factor
$\exp\!\bigl(i\mu_p v_i^{(\alpha)} \ell_0/\hbar_{\mathrm{eff}}\bigr)$
is a discrete analogue of the de~Broglie phase $e^{ip\ell/\hbar}$:
it encodes the kinematic momentum of particle $i$ in the $\alpha$-direction into the
phase of the $\alpha$-th color component.
Nodes moving in different directions therefore carry genuinely different
internal phases, even if they experience the same force magnitude.
This momentum-phase coupling ensures that the resulting parallel transport
$U_3(e)$ depends not only on the local force profile but also on how the color
triplet changes from one site to the next along the oriented edge.

The regularization $\epsilon_c>0$ ensures that $\mathbf c_x$ is well-defined
even when the viscous force vanishes, e.g.\ for an isolated particle or when
$\nu_{\mathrm{visc}}=0$.
If $\mathbf c_x=0$, then
\[
\Re\!\bigl(\mathbf c_x^\dagger\lambda_a(\mathbf c_y-\mathbf c_x)\bigr)=0
\]
for all outgoing links from $x$, hence $U_3(x\to y)=I$ for all such links.
Thus nodes with trivial source color couple trivially through this link
transport. When $\mathbf c_y=0$ and $\mathbf c_x\neq0$, the link need not be
trivial: it records the local decay of the color field along the edge.

The normalization satisfies $\|\mathbf c_x\|^2<1$ for all $x$ (Lemma~\ref{lem:color-norm}).
Combined with $|\lambda_a|_\mathrm{op}\le 2$ for all Gell-Mann generators,
this gives the operator bound
\[
\|\Theta_e\|_{\mathrm{op}}
\le \frac{g_3}{2}\sum_{a=1}^8
\bigl|\Re\!\bigl(\mathbf c_x^\dagger\lambda_a(\mathbf c_y-\mathbf c_x)\bigr)\bigr|
\|\lambda_a\|_{\mathrm{op}}
\le \frac{g_3}{2}\cdot 8\cdot \|\mathbf c_x\|\,\|\mathbf c_y-\mathbf c_x\|\cdot 2
< 16g_3,
\]
so the gauge coupling strength is controlled uniformly over all links and runs.
\end{remark}

\begin{lemma}[Normalization of color vectors]
\label{lem:color-norm}
For every site $x\in\mcN$,
\[
\mathbf c_x^\dagger\mathbf c_x
=\frac{\tilde{\mathbf c}_x^\dagger\tilde{\mathbf c}_x}{\tilde{\mathbf c}_x^\dagger\tilde{\mathbf c}_x+\epsilon_c^2}\in[0,1),
\]
so $\mathbf c_x$ is well-defined and $\|\mathbf c_x\|<1$.
\end{lemma}
\begin{proof}
Since $\epsilon_c>0$, the denominator $\tilde{\mathbf c}_x^\dagger\tilde{\mathbf c}_x+\epsilon_c^2$
is strictly positive.
The ratio $r/(r+\epsilon_c^2)$ for $r=\tilde{\mathbf c}_x^\dagger\tilde{\mathbf c}_x\ge0$ lies in $[0,1)$,
with the upper bound strict because $\epsilon_c^2>0$.
\end{proof}

\begin{lemma}[Color reparametrization]
\label{lem:color-reparam}
Let $\Omega:\mcN\to SU(3)$ be an arbitrary map.
Define $\mathbf c_x^\Omega:=\Omega(x)\mathbf c_x$ for each $x\in\mcN$.
Then $(\mathbf c_x^\Omega)^\dagger\mathbf c_x^\Omega=\mathbf c_x^\dagger\mathbf c_x$
and $\|\mathbf c_x^\Omega\|=\|\mathbf c_x\|$.
\end{lemma}
\begin{proof}
Since $\Omega(x)\in U(3)$, $\Omega(x)^\dagger\Omega(x)=I$.
Hence
\[
(\mathbf c_x^\Omega)^\dagger\mathbf c_x^\Omega
=\mathbf c_x^\dagger\Omega(x)^\dagger\Omega(x)\mathbf c_x=\mathbf c_x^\dagger\mathbf c_x.
\]
\end{proof}

\begin{definition}[Gell-Mann generators]
\label{def:gellmann}
The eight Gell-Mann matrices $\lambda_a\in\C^{3\times3}$, $a=1,\dots,8$, are
\[
\lambda_1=\begin{pmatrix}0&1&0\\1&0&0\\0&0&0\end{pmatrix},\;
\lambda_2=\begin{pmatrix}0&-i&0\\ i&0&0\\0&0&0\end{pmatrix},\;
\lambda_3=\begin{pmatrix}1&0&0\\0&-1&0\\0&0&0\end{pmatrix},
\]
\[
\lambda_4=\begin{pmatrix}0&0&1\\0&0&0\\1&0&0\end{pmatrix},\;
\lambda_5=\begin{pmatrix}0&0&-i\\0&0&0\\ i&0&0\end{pmatrix},\;
\lambda_6=\begin{pmatrix}0&0&0\\0&0&1\\0&1&0\end{pmatrix},
\]
\[
\lambda_7=\begin{pmatrix}0&0&0\\0&0&-i\\0&i&0\end{pmatrix},\;
\lambda_8=\frac1{\sqrt3}\begin{pmatrix}1&0&0\\0&1&0\\0&0&-2\end{pmatrix}.
\]
Each $\lambda_a$ is Hermitian and $\operatorname{tr}\lambda_a=0$.
The generators $\{i\lambda_a/2\}_{a=1}^8$ form a basis of the Lie algebra
$\mathfrak{su}(3)$ of traceless skew-Hermitian matrices.
\end{definition}

\begin{lemma}[Gell-Mann orthogonality]
\label{lem:gellmann}
The Gell-Mann matrices satisfy
\[
\tr(\lambda_a\lambda_b)=2\delta_{ab},\qquad a,b\in\{1,\dots,8\}.
\]
\end{lemma}
\begin{proof}
Direct matrix multiplication: for $a=b$ each diagonal entry of $\lambda_a^2$
contributes, giving trace $2$; for $a\ne b$ the off-diagonal structure ensures
the trace vanishes.
This is a finite componentwise check whose result is standard.
\end{proof}

\begin{lemma}[Hermitian traceless span]
\label{lem:hermitian-span}
If $\alpha_1,\dots,\alpha_8\in\R$, then $H:=\sum_{a=1}^8\alpha_a\lambda_a$
is Hermitian and traceless.
\end{lemma}
\begin{proof}
Each $\lambda_a$ is Hermitian and traceless by Definition~\ref{def:gellmann}.
Real linear combinations preserve both properties.
\end{proof}

\begin{definition}[Generator and transport on links]
\label{def:transport}
For each oriented link $e=(x\to y)\in\mcE$ define the algebra element
\[
\mathcal A_e:=\frac12\sum_{a=1}^8
\Re\!\bigl(\mathbf c_x^\dagger\lambda_a(\mathbf c_y-\mathbf c_x)\bigr)\lambda_a
\in\{H\in\C^{3\times3}:H^\dagger=H,\ \tr H=0\},
\qquad
\Theta_e:=g_3\,\mathcal A_e.
\]
and the $SU(3)$ link transport
\[
U_3(e):=\exp(i\Theta_e)\in SU(3),
\qquad U_3(e^{-1}):=U_3(e)^{-1}.
\]
\end{definition}

\begin{remark}[Justification of the transport construction]
\label{rem:transport-justification}
The formula
\[
\mathcal A_e=\frac12\sum_a
\Re\!\bigl(\mathbf c_x^\dagger\lambda_a(\mathbf c_y-\mathbf c_x)\bigr)\lambda_a
\]
is the orthogonal projection of the discrete color increment
$\mathbf c_y-\mathbf c_x$ against the source color vector onto the Hermitian
traceless subspace spanned by $\{\lambda_a\}_{a=1}^8$.
To see this, recall that the Hermitian traceless matrices carry the inner
product $\langle A,B\rangle=\frac{1}{2}\tr(AB)$, under which $\{\lambda_a\}$ is
an orthogonal basis with $\langle\lambda_a,\lambda_b\rangle=\delta_{ab}$ by
Lemma~\ref{lem:gellmann}.
The $a$-th component of the projected discrete derivative is
\[
\bigl\langle \lambda_a,\tfrac12\bigl(\mathbf c_x(\mathbf c_y-\mathbf c_x)^\dagger
+(\mathbf c_y-\mathbf c_x)\mathbf c_x^\dagger\bigr)\bigr\rangle
=\frac12\tr\!\Bigl(\lambda_a\bigl(\mathbf c_x(\mathbf c_y-\mathbf c_x)^\dagger
+(\mathbf c_y-\mathbf c_x)\mathbf c_x^\dagger\bigr)\Bigr)
=\Re\!\bigl(\mathbf c_x^\dagger\lambda_a(\mathbf c_y-\mathbf c_x)\bigr),
\]
where we used $\lambda_a^\dagger=\lambda_a$.
Since the resulting coefficients are real, $\mathcal A_e$ is exactly this
projection. The coupling is $\Theta_e=g_3\mathcal A_e$.

The factor $g_3/2$ is chosen to match the continuum covariant derivative
$D_\mu=\partial_\mu-ig_3 A_\mu^a\lambda_a/2$:
if the color field varies smoothly and the edge displacement is small, then
$\mathbf c_y-\mathbf c_x$ is first order in the edge length, so the transport
linearizes to
\[
U_3(e)=I+i\Theta_e+O(g_3^2)
=I+ig_3\,\mathcal A_e+O(g_3^2)
=I+\frac{ig_3}{2}\sum_a
\Re\!\bigl(\mathbf c_x^\dagger\lambda_a(\mathbf c_y-\mathbf c_x)\bigr)\lambda_a
+O(g_3^2),
\]
which is the first-order approximation to a continuum parallel transporter
$P\exp(ig_3\int A_\mu\,dx^\mu)$.
\end{remark}

\begin{proposition}[Link transport is $SU(3)$]
\label{prop:su3-link}
For every $e\in\mcE$, $U_3(e)\in SU(3)$ and $U_3(e^{-1})=U_3(e)^{-1}$.
\end{proposition}
\begin{proof}
Each coefficient $\Re(\mathbf c_x^\dagger\lambda_a\mathbf c_y)$ is real.
By Lemma~\ref{lem:hermitian-span}, $\Theta_e$ is Hermitian and traceless, so
$i\Theta_e$ is skew-Hermitian and traceless.
Hence $U_3(e)=\exp(i\Theta_e)$ is unitary.
Since $\tr\Theta_e=0$, the identity $\det\exp(M)=\exp(\tr M)$ gives
$\det U_3(e)=\exp(i\tr\Theta_e)=1$.
Therefore $U_3(e)\in SU(3)$.
The inverse relation $U_3(e^{-1})=U_3(e)^{-1}$ is imposed by convention.
\end{proof}

\begin{remark}
There are no oriented self-loops in $\mcE$ because $(i,i)\notin\mathcal P_t$,
so source and target are distinct by definition.
If one formally assigns the link transport formula to a self-loop, it still lands in
$SU(3)$ because $\mathcal A_e$ is Hermitian and traceless.
\end{remark}

\section{Holonomy and gauge invariance}
\label{sec:holonomy}

\begin{definition}[Path holonomy]
\label{def:path}
For an oriented path
$\gamma=(v_0\to v_1,\dots,v_{m-1}\to v_m)$ with $v_k\in\mcN$,
define the ordered product (left-to-right composition)
\[
U_3(\gamma):=U_3(v_0\to v_1)\,U_3(v_1\to v_2)\cdots U_3(v_{m-1}\to v_m)
=\prod_{k=1}^m U_3(v_{k-1}\to v_k).
\]
This ordering convention uses the usual link ordering in lattice gauge theory,
where $U_\mu(x)$ transports from site $x$ to site $x+\hat\mu$.
If $v_m=v_0$, $\gamma$ is a loop and we define the Wilson loop trace
$W_3(\gamma):=\tr U_3(\gamma)$.
\end{definition}

\begin{definition}[Gauge map]
\label{def:gauge}
A gauge map is any map $\Omega:\mcN\to SU(3)$.
For a link $e=(x\to y)\in\mcE$ define the gauge-transformed transport
\[
U_3^\Omega(e):=\Omega(y)\,U_3(e)\,\Omega(x)^{-1}.
\]
Transformed color vectors are $\mathbf c_x^\Omega:=\Omega(x)\mathbf c_x$.
\end{definition}

\begin{proposition}[Gauge invariance of Wilson traces]
\label{prop:gauge-invariance}
For any closed loop $\gamma=(v_0,\dots,v_m=v_0)$ and any gauge map $\Omega$,
\[
\tr\left(U_3^\Omega(\gamma)\right)=\tr\left(U_3(\gamma)\right).
\]
\end{proposition}
\begin{proof}
Compute the gauge-transformed holonomy:
\[
U_3^\Omega(\gamma)
=\prod_{k=1}^m \Omega(v_k)\,U_3(v_{k-1}\to v_k)\,\Omega(v_{k-1})^{-1}.
\]
Consecutive factors telescope: $\Omega(v_k)$ from step $k$ cancels
$\Omega(v_k)^{-1}$ from step $k+1$, since $v_k$ appears as the target of step $k$
and the source of step $k+1$.
After all $m$ steps, the remaining factors are $\Omega(v_m)=\Omega(v_0)$ on the left
and $\Omega(v_0)^{-1}$ on the right, giving
\[
U_3^\Omega(\gamma)=\Omega(v_0)\,U_3(\gamma)\,\Omega(v_0)^{-1}.
\]
Trace cyclicity then gives
$\tr(U_3^\Omega(\gamma))=\tr(\Omega(v_0)^{-1}\Omega(v_0)\,U_3(\gamma))=\tr(U_3(\gamma))$.
\end{proof}

\begin{definition}[Canonical elementary plaquettes]
\label{def:plaquettes}
For each time slice $0\le t<T$ and each ordered pair
$(i,j)\in\mathcal P_t\cap\mathcal P_{t+1}$, define the elementary plaquette
\[
P_{ij,t}:=(n_{i,t}\to n_{i,t+1}\to n_{j,t+1}\to n_{j,t}\to n_{i,t}).
\]
Its boundary consists of one forward temporal link for particle $i$, one forward
interaction link at time $t+1$, one backward temporal link for particle $j$,
and one backward interaction link at time $t$.
Set
\[
\mathcal C_\square:=\{P_{ij,t}:\ (i,j)\in\mathcal P_t\cap\mathcal P_{t+1},\ 0\le t<T\}.
\]
\end{definition}

\begin{definition}[Plaquette geometry]
\label{def:plaquette-geometry}
For $P=P_{ij,t}\in\mathcal C_\square$, define the base point
\[
z_P:=\Xi(n_{i,t})\in\R^4,
\]
the temporal and interaction edge displacements at the base point
\[
\tau_{i,t}:=\Xi(n_{i,t+1})-\Xi(n_{i,t})\in\R^4,
\qquad
\sigma_{ij,t}:=\Xi(n_{j,t})-\Xi(n_{i,t})\in\R^4,
\]
and the oriented area bivector
\[
\Sigma_P^{\mu\nu}:=\tau_{i,t}^\mu\sigma_{ij,t}^\nu-\tau_{i,t}^\nu\sigma_{ij,t}^\mu.
\]
Its area is
\[
A_P:=|\Sigma_P|:=\left(\frac12\sum_{\mu,\nu=0}^3
\Sigma_P^{\mu\nu}\Sigma_P^{\mu\nu}\right)^{1/2}.
\]
\end{definition}

\begin{definition}[Plaquette holonomy and discrete field strength]
\label{def:plaquette-field-strength}
For $P=P_{ij,t}\in\mathcal C_\square$, define the plaquette holonomy by the
ordered product
\[
U_3(P):=
U_3(n_{i,t}\to n_{i,t+1})\,
U_3(n_{i,t+1}\to n_{j,t+1})\,
U_3(n_{j,t+1}\to n_{j,t})\,
U_3(n_{j,t}\to n_{i,t}).
\]
Whenever $A_P>0$ and $U_3(P)$ lies in the principal logarithm domain, define the
discrete plaquette field strength by
\[
F_P:=\frac{1}{i g_3 A_P}\log U_3(P)\in\mathfrak{su}(3).
\]
\end{definition}

\begin{definition}[Wilson action]
\label{def:action}
Set $\beta:=6/g_3^2$.
The Wilson plaquette action is
\[
S_{\mathrm W}:=\beta\sum_{P\in\mathcal C_\square}\left(1-\frac13\Re\tr U_3(P)\right).
\]
\end{definition}

\begin{remark}
The normalization $1-\frac{1}{3}\Re\tr U_3(P)$ vanishes when $U_3(P)=I$ and is
non-negative for all $U_3(P)\in SU(3)$, since
$|\tr U|/3\le1$ by the triangle inequality applied to the three eigenvalues of $U$.
In the small-loop regime, Theorem~\ref{thm:small-loop-expansion} identifies the
leading-order term with a Yang--Mills curvature density.
\end{remark}

\begin{proposition}[Gauge invariance of Wilson action]
\label{prop:wilson-gauge}
For any gauge map $\Omega$, $S_{\mathrm W}^\Omega=S_{\mathrm W}$.
\end{proposition}
\begin{proof}
Each loop trace $\Re\tr U_3(C)$ is invariant under gauge transformation by
Proposition~\ref{prop:gauge-invariance}.
The sum and the constant term are unchanged, so $S_{\mathrm W}^\Omega=S_{\mathrm W}$.
\end{proof}

\section{Matter sector on the same lattice}
\label{sec:matter}

\begin{definition}[Matter fields and actions]
\label{def:matter}
For each site $x\in\mcN$ choose a fermionic spinor $\psi_x\in\C^3$,
its conjugate $\bar\psi_x\in(\C^3)^\dagger$, and a real scalar $\phi_x\in\R$.
Fix $m_F>0,\ m_\phi>0$, and a local potential $V_\phi:\R\to\R$.
For $e=(x\to y)\in\mcE$ set
\[
\ell_{xy}:=\|\Xi(y)-\Xi(x)\|+\eta=\ell_e+\eta
\]
with $\eta>0$.
The covariant fermionic action is
\[
S_{\mathrm F}:=\sum_{e\in\mcE}\bar\psi_{t(e)}\left(\psi_{t(e)}-U_3(e)\psi_{s(e)}\right)
+m_F\sum_{x\in\mcN}\bar\psi_x\psi_x,
\]
and the covariant scalar action is
\[
S_\phi:=\frac12\sum_{e\in\mcE}\frac{(\phi_{t(e)}-\phi_{s(e)})^2}{\ell_{xy}}
+\sum_{x\in\mcN}\Bigl(\frac{m_\phi^2}{2}\phi_x^2+V_\phi(\phi_x)\Bigr).
\]
The gauge group acts on matter fields by
\[
\psi_x\mapsto\Omega(x)\psi_x,\qquad
\bar\psi_x\mapsto\bar\psi_x\Omega(x)^{-1},\qquad
\phi_x\mapsto\phi_x.
\]
\end{definition}

\begin{remark}[Distance weighting in the scalar action]
\label{rem:distance-weight}
The factor $1/\ell_{xy}$ in the scalar kinetic term is a discrete analogue of
the continuum normalization $(\nabla_\mu\phi)^2$.
In the continuum, the square of a derivative has dimensions $[\phi]^2/[\text{length}]^2$,
so a finite difference $(\phi_y-\phi_x)^2$ must be divided by $\|\mathbf x(y)-\mathbf x(x)\|$
(one power of length rather than two) for a first-order finite-difference scheme.
In the spacetime lattice used here, the relevant finite-difference length is the
four-dimensional edge norm $\|\Xi(y)-\Xi(x)\|$.
Specifically, if $\phi_{t(e)}-\phi_{s(e)}\approx (\partial_\mu\phi)\,\xi_e^\mu$
with $\|\Xi(y)-\Xi(x)\|\approx\ell_{xy}$, then
$(\phi_{t(e)}-\phi_{s(e)})^2/\ell_{xy}\approx \ell_{xy}\,|\nabla\phi|^2$,
and summing over a regular lattice with spacing $a=\ell_{xy}$ recovers
$a^d\,|\nabla\phi|^2$ as in the usual scalar lattice action.

The parameter $\eta>0$ prevents division by zero for coincident site positions
(self-loops, or particles occupying the same spacetime point).
In practice $\eta\ll\ell_{\mathrm{visc}}$ so that $\ell_{xy}\approx\|\Xi(y)-\Xi(x)\|$
for all non-coincident links.
\end{remark}

\begin{proposition}[Matter gauge invariance]
\label{prop:matter-gauge}
Under the gauge action of Definition~\ref{def:matter},
\[
S_{\mathrm F}^\Omega=S_{\mathrm F},\qquad S_\phi^\Omega=S_\phi.
\]
\end{proposition}
\begin{proof}
For each fermionic link term,
\begin{align*}
\bar\psi_{t(e)}\!\left(\psi_{t(e)}-U_3(e)\psi_{s(e)}\right)
&\mapsto
\bar\psi_{t(e)}\Omega(t(e))^{-1}
\!\left(\Omega(t(e))\psi_{t(e)}-\Omega(t(e))U_3(e)\Omega(s(e))^{-1}\Omega(s(e))\psi_{s(e)}\right)\\
&=\bar\psi_{t(e)}\!\left(\psi_{t(e)}-U_3(e)\psi_{s(e)}\right),
\end{align*}
where we used $\Omega(t(e))^{-1}\Omega(t(e))=I$ and $\Omega(s(e))^{-1}\Omega(s(e))=I$.
The mass term $m_F\bar\psi_x\psi_x$ is invariant because $\Omega(x)^{-1}\Omega(x)=I$ pointwise.
The scalar action $S_\phi$ depends only on scalar fields, which are gauge-invariant,
so $S_\phi^\Omega=S_\phi$.
\end{proof}

\begin{remark}[Matter sector on the lattice]
The same finite link set $\mcE$ supports both matter sectors.
The fermionic term couples oriented transport through $U_3(e)$, while the
scalar term uses the same link differences as a discrete kinetic energy.
\end{remark}

\section{Finite lattice-QFT data package}
\label{sec:package}

\begin{theorem}[Finite $SU(3)$ data package]
\label{thm:package}
Given a run $\mathscr R$ as in Definition~\ref{def:run} and the constructions
in Sections~\ref{sec:run}--\ref{sec:matter}, the tuple
\[
\mathcal Q=
\Bigl(
T,\Delta t,\{\mathcal A(t)\},\{\mathcal P_t\},\mcN,E_{\mathrm{int}},E_{\mathrm{time}},\mcE,\mathcal C_\square,
\{x_i(t),v_i(t),\Phi_i(t)\}_{i,t},K_\rho,\nu_{\mathrm{visc}},
\hbar_{\mathrm{eff}},\mu_p,\ell_0,\epsilon_c,g_3,\beta,\eta,
\{\Xi(x)\}_{x\in\mcN},\{\Sigma_P,A_P,U_3(P)\}_{P\in\mathcal C_\square},
\{U_3(e)\}_{e\in\mcE},S_{\mathrm W},S_{\mathrm F},S_\phi
\Bigr)
\]
is finite and explicit.
It carries:
\begin{enumerate}[leftmargin=1.5em]
\item a finite spacetime lattice $(\mcN,\mcE)$ with interaction and temporal links;
\item a finite canonical plaquette family $\mathcal C_\square$ for the Wilson action;
\item explicit plaquette geometry $(z_P,\Sigma_P,A_P)$ and plaquette holonomies $U_3(P)$;
\item $SU(3)$ parallel transport maps $\{U_3(e)\}_{e\in\mcE}$;
\item gauge-invariant Wilson plaquette traces and Wilson action $S_{\mathrm W}$;
\item gauge-covariant fermionic and scalar matter actions $S_{\mathrm F}$, $S_\phi$.
\end{enumerate}
All objects required for a subsequent axiom-verification paper are present.
\end{theorem}
\begin{proof}
Finiteness of $(\mcN,\mcE)$ is Proposition~\ref{prop:finiteness}.
Finiteness of $\mathcal C_\square$ follows from Definition~\ref{def:plaquettes}
because it is indexed by the finite set of triples $(i,j,t)$ with
$0\le t<T$ and $(i,j)\in\mathcal P_t\cap\mathcal P_{t+1}$.
The spacetime embedding, area bivectors, and plaquette holonomies are finite
families because they are explicit algebraic functions of $\mcN$ and
$\mathcal C_\square$.
Finiteness of $\{U_3(e)\}$ follows because $\mcE$ is finite and each $U_3(e)$
is defined by a finite matrix exponential.
The actions $S_{\mathrm W}$, $S_{\mathrm F}$, $S_\phi$ are finite sums over
finite index sets by the same finiteness.
Gauge invariance of $S_{\mathrm W}$ is Proposition~\ref{prop:wilson-gauge};
gauge invariance of $S_{\mathrm F}$ and $S_\phi$ is
Proposition~\ref{prop:matter-gauge}.
\end{proof}

\section{Discussion}
\label{sec:discussion}

\subsection{Lattice-gauge interpretation}

The lattice $(\mcN,\mcE)$ is finite and data dependent.
The site set is a finite subset of spacetime determined by the particle data,
while the link set contains both temporal worldline links and same-time
interaction links.
The link variables $U_3(e)$ are not independent: they are
\emph{derived} from the trajectory data via the color construction of
Definitions~\ref{def:colors}--\ref{def:transport}.
The gauge group, the coupling strength $g_3$, and the parallel transport
all emerge from the viscous dynamics rather than being postulated.

The Wilson action is written on the canonical plaquette family
$\mathcal C_\square$ from Definition~\ref{def:plaquettes}, and the matter
actions use the same spacetime links.
All gauge-sector observables are therefore evaluated on explicit elementary
2-cells rather than on an arbitrary loop family.

\subsection{Thermodynamic limit of graph observables}

The convergence paper proves exchangeability of the invariant laws, a finite
de~Finetti representation for the empirical measure, collapse of the mixing
measure to the unique stationary mean-field law $\pi_\infty$, and the same
$N$-uniform logarithmic Sobolev constant in the limit~\cite{convergence}.
These results are exactly the input needed to pass the discrete interaction
graph itself to the stationary mean-field regime.

\begin{definition}[Pairwise link rule and empirical edge measure]
\label{def:edge-measure}
Fix a recorded time slice $t\in\Tset$ and, for the $N$-particle system, write
\[
Z_i^N(t):=(X_i^N(t),V_i^N(t))\in\R^6,
\qquad
L_t^N:=\frac1N\sum_{i=1}^N \delta_{Z_i^N(t)}.
\]
Assume that the same-time link set $\mathcal P_t^N$ is generated by the same
label-invariant pair rule for every $N$, in the sense that there exists a
bounded measurable indicator
\[
\chi_t:\R^6\times\R^6\to\{0,1\}
\]
such that
\[
\mathbf 1_{\{(i,j)\in\mathcal P_t^N\}}
=
\chi_t(Z_i^N(t),Z_j^N(t))
\qquad
(i\ne j).
\]
Define the empirical edge measure by
\[
\Gamma_t^N
:=
\frac1{N^2}\sum_{i\ne j}
\mathbf 1_{\{(i,j)\in\mathcal P_t^N\}}
\delta_{(Z_i^N(t),Z_j^N(t))}.
\]
\end{definition}

\begin{theorem}[Stationary thermodynamic limit of the interaction graph]
\label{thm:edge-measure-limit}
Assume the stationary thermodynamic limit from the companion convergence paper:
under the invariant law of the $N$-particle system, the random empirical
measure $L_t^N$ converges in law to the deterministic stationary mean-field law
$\pi_\infty$ as $N\to\infty$.
Let $\Phi:\R^6\times\R^6\to\R$ be bounded and continuous, and assume that
$\chi_t\Phi$ is continuous at $(\pi_\infty\otimes\pi_\infty)$-almost every pair.
Then
\[
\int_{\R^6\times\R^6}\Phi\,d\Gamma_t^N
\xrightarrow[N\to\infty]{\mathbb P}
\iint_{\R^6\times\R^6}
\chi_t(z,z')\,\Phi(z,z')\,\pi_\infty(dz)\,\pi_\infty(dz').
\]
Equivalently,
\[
\frac1{N^2}\sum_{(i,j)\in\mathcal P_t^N}\Phi(Z_i^N(t),Z_j^N(t))
\xrightarrow[N\to\infty]{\mathbb P}
\iint_{\R^6\times\R^6}
\chi_t(z,z')\,\Phi(z,z')\,\pi_\infty(dz)\,\pi_\infty(dz').
\]
If the number of recorded slices $T+1$ is fixed, then the same convergence
holds after summing over $t\in\Tset$.
\end{theorem}
\begin{proof}
Write
\[
\Psi_t(z,z'):=\chi_t(z,z')\,\Phi(z,z').
\]
By Definition~\ref{def:edge-measure},
\begin{align*}
\int \Phi\,d\Gamma_t^N
&=
\frac1{N^2}\sum_{i\ne j}\Psi_t(Z_i^N(t),Z_j^N(t))\\
&=
\frac1{N^2}\sum_{i,j}\Psi_t(Z_i^N(t),Z_j^N(t))
-\frac1{N^2}\sum_{i=1}^N \Psi_t(Z_i^N(t),Z_i^N(t))\\
&=
\iint_{\R^6\times\R^6}\Psi_t(z,z')\,L_t^N(dz)\,L_t^N(dz')
-R_t^N,
\end{align*}
where
\[
|R_t^N|\le \frac{\|\Psi_t\|_\infty}{N}.
\]
Hence $R_t^N\to0$ in probability.

Because $L_t^N$ converges in law to the deterministic measure $\pi_\infty$,
the product measure $L_t^N\otimes L_t^N$ converges in law to
$\pi_\infty\otimes\pi_\infty$ by the continuous mapping theorem.
The boundedness and almost-everywhere continuity assumption on $\Psi_t$ then give
\[
\iint \Psi_t(z,z')\,L_t^N(dz)\,L_t^N(dz')
\xrightarrow[N\to\infty]{\mathbb P}
\iint \Psi_t(z,z')\,\pi_\infty(dz)\,\pi_\infty(dz').
\]
Combining the last two displays proves the first claim.
The displayed sum over $(i,j)\in\mathcal P_t^N$ is exactly the same quantity
written in index notation, so the equivalent form follows.
If $T+1$ is fixed, summing the convergences over finitely many times preserves
convergence in probability.
\end{proof}

\begin{corollary}[Thermodynamic limit of weighted link sums]
\label{cor:weighted-edge-limit}
Let
\[
W_t:\R^6\times\R^6\times \mathcal P(\R^6)\to\R
\]
be bounded, and assume that $W_t$ is jointly continuous at every point
$(z,z',\pi_\infty)$ for which $\chi_t(z,z')=1$.
Then for every bounded continuous $\Phi$,
\[
\frac1{N^2}\sum_{(i,j)\in\mathcal P_t^N}
W_t(Z_i^N(t),Z_j^N(t);L_t^N)\,\Phi(Z_i^N(t),Z_j^N(t))
\xrightarrow[N\to\infty]{\mathbb P}
\iint
\chi_t(z,z')\,W_t(z,z';\pi_\infty)\,\Phi(z,z')\,
\pi_\infty(dz)\,\pi_\infty(dz').
\]
\end{corollary}
\begin{proof}
Set
\[
\Psi_t^N(z,z')
:=
\chi_t(z,z')\,W_t(z,z';L_t^N)\,\Phi(z,z'),
\qquad
\Psi_t^\infty(z,z')
:=
\chi_t(z,z')\,W_t(z,z';\pi_\infty)\,\Phi(z,z').
\]
Since $L_t^N\to\pi_\infty$ in law and the limit is deterministic,
$W_t(z,z';L_t^N)\to W_t(z,z';\pi_\infty)$ in probability for each continuity
point $(z,z',\pi_\infty)$.
Boundedness then lets us pass to the limit in the same decomposition used in
Theorem~\ref{thm:edge-measure-limit}, with $\Psi_t^N$ in place of $\Psi_t$.
The diagonal contribution is still $O(N^{-1})$, so the result follows.
\end{proof}

\begin{remark}
Theorem~\ref{thm:edge-measure-limit} is the thermodynamic step needed to replace
discrete interaction-link sums by mean-field integrals without introducing any
external discretization of the limiting law.
To pass the Wilson action itself to the thermodynamic limit one also needs a
plaquette-level statement, because the canonical plaquettes depend on two
adjacent time slices.
\end{remark}

\begin{definition}[Adjacent-time empirical path and plaquette measures]
\label{def:plaquette-measure}
Fix $0\le t<T$ and write
\[
\widetilde Z_i^N(t):=\bigl(Z_i^N(t),Z_i^N(t+1)\bigr)\in\R^{12},
\qquad
M_t^N:=\frac1N\sum_{i=1}^N \delta_{\widetilde Z_i^N(t)}.
\]
Define the persistent-plaquette indicator by
\[
\chi_{\square,t}\bigl((z_0,z_1),(w_0,w_1)\bigr):=\chi_t(z_0,w_0)\,\chi_{t+1}(z_1,w_1),
\]
and the empirical plaquette measure by
\[
\Gamma_{\square,t}^N
:=
\frac1{N^2}\sum_{i\ne j}
\chi_{\square,t}\bigl(\widetilde Z_i^N(t),\widetilde Z_j^N(t)\bigr)\,
\delta_{\bigl(\widetilde Z_i^N(t),\widetilde Z_j^N(t)\bigr)}.
\]
\end{definition}

\begin{theorem}[Stationary thermodynamic limit of plaquette observables]
\label{thm:plaquette-measure-limit}
Assume that, under the invariant law of the $N$-particle system, the adjacent-time
empirical path measure $M_t^N$ converges in law to a deterministic mean-field
path law $\Lambda_t^\infty$ on $\R^{12}$ as $N\to\infty$.
Let $\Phi:(\R^{12})^2\to\R$ be bounded and continuous, and assume that
$\chi_{\square,t}\Phi$ is continuous at
$(\Lambda_t^\infty\otimes\Lambda_t^\infty)$-almost every pair.
Then
\[
\int_{(\R^{12})^2}\Phi\,d\Gamma_{\square,t}^N
\xrightarrow[N\to\infty]{\mathbb P}
\iint_{(\R^{12})^2}
\chi_{\square,t}(\zeta,\zeta')\,\Phi(\zeta,\zeta')\,
\Lambda_t^\infty(d\zeta)\,\Lambda_t^\infty(d\zeta').
\]
\end{theorem}
\begin{proof}
Set
\[
\Psi_{\square,t}(\zeta,\zeta'):=\chi_{\square,t}(\zeta,\zeta')\,\Phi(\zeta,\zeta').
\]
By Definition~\ref{def:plaquette-measure},
\begin{align*}
\int \Phi\,d\Gamma_{\square,t}^N
&=
\frac1{N^2}\sum_{i\ne j}
\Psi_{\square,t}\bigl(\widetilde Z_i^N(t),\widetilde Z_j^N(t)\bigr)\\
&=
\iint \Psi_{\square,t}(\zeta,\zeta')\,M_t^N(d\zeta)\,M_t^N(d\zeta')
\;-\;R_{\square,t}^N,
\end{align*}
where $|R_{\square,t}^N|\le \|\Psi_{\square,t}\|_\infty/N$.
Hence $R_{\square,t}^N\to0$ in probability.
Because $M_t^N$ converges in law to the deterministic measure $\Lambda_t^\infty$,
the product measure $M_t^N\otimes M_t^N$ converges in law to
$\Lambda_t^\infty\otimes\Lambda_t^\infty$.
The boundedness and almost-everywhere continuity of $\Psi_{\square,t}$ then yield
\[
\iint \Psi_{\square,t}(\zeta,\zeta')\,M_t^N(d\zeta)\,M_t^N(d\zeta')
\xrightarrow[N\to\infty]{\mathbb P}
\iint \Psi_{\square,t}(\zeta,\zeta')\,
\Lambda_t^\infty(d\zeta)\,\Lambda_t^\infty(d\zeta').
\]
Combining the last two displays proves the theorem.
\end{proof}

\begin{remark}
Theorems~\ref{thm:edge-measure-limit} and~\ref{thm:plaquette-measure-limit}
are the thermodynamic statements needed for the gauge sector.
The first passes link observables to the mean-field limit; the second does the
same for the canonical plaquettes used by the Wilson action.
No external discretization of the limiting law is introduced: the dynamic graph
and its elementary 2-cells are carried directly to the mean-field regime by the
same label-invariant graph rule.
\end{remark}

\subsection{Continuum limit}

The thermodynamic limit identifies the mean-field objects associated with links
and plaquettes.
To recover a continuum gauge theory from the same data one needs a second limit
in which the mesh size tends to zero, the color field becomes smooth, and the
plaquette family samples the ambient spacetime geometry in an asymptotically
isotropic way.
The following definition isolates exactly those hypotheses.

\begin{definition}[Admissible mean-field refinement regime]
\label{def:refinement}
Consider a sequence of data sets
\[
\mathscr R^{(N)}=
\bigl(\mcN^{(N)},\mcE^{(N)},\mathcal C_\square^{(N)},\{U_3^{(N)}(e)\}\bigr),
\qquad N\to\infty.
\]
Write
\[
a_N:=\max_{e\in\mcE^{(N)}} \ell_e^{(N)}.
\]
The sequence is an \emph{admissible mean-field refinement regime} if the following hold.
\begin{enumerate}[leftmargin=1.5em]
\item \textbf{Vanishing mesh.}
$a_N\to0$ and $\Delta t_N\le C a_N$.
\item \textbf{Shape-regular plaquettes.}
There exist constants $0<c_-<c_+<\infty$ and $C_{\mathrm{reg}}<\infty$ such that
for every $P\in\mathcal C_\square^{(N)}$,
\[
c_- a_N^2\le A_P^{(N)}\le c_+ a_N^2,
\]
and if $P=P_{ij,t}$, then
\[
\|\tau_{j,t}^{(N)}-\tau_{i,t}^{(N)}\|
+\|\sigma_{ij,t+1}^{(N)}-\sigma_{ij,t}^{(N)}\|
\le C_{\mathrm{reg}}\,a_N^2.
\]
\item \textbf{Smooth color-field limit.}
There exists a smooth embedded four-dimensional manifold $M\subset\R^4$ with
Riemannian metric $g$ and a $C^2$ color field
\[
\mathbf c:M\to\C^3
\]
such that, after identifying $\Xi^{(N)}(\mcN^{(N)})\subset M$,
\[
\sup_{x\in\mcN^{(N)}}\|\mathbf c_x^{(N)}-\mathbf c(\Xi^{(N)}(x))\|
\le C_{\mathbf c}\,a_N^2.
\]
\item \textbf{Bivector density convergence.}
For every compactly supported continuous scalar test function $f$ on $M$,
\[
\sum_{P\in\mathcal C_\square^{(N)}}
\Sigma_P^{(N),\mu\nu}\Sigma_P^{(N),\rho\sigma}\,f(z_P^{(N)})
\longrightarrow
\frac18\int_M
\bigl(g^{\mu\rho}g^{\nu\sigma}-g^{\mu\sigma}g^{\nu\rho}\bigr)
f(z)\,d\mathrm{vol}_g(z).
\]
\end{enumerate}
\end{definition}

\begin{theorem}[Short-link consistency and continuum connection]
\label{thm:short-link-connection}
Assume the admissible refinement regime of Definition~\ref{def:refinement}.
Define the Hermitian $\mathfrak{su}(3)$-valued connection on $M$ by
\[
A_\mu(z):=\frac12\sum_{a=1}^8
\Re\!\bigl(\mathbf c(z)^\dagger\lambda_a\,\partial_\mu\mathbf c(z)\bigr)\lambda_a,
\qquad z\in M.
\]
Then for every oriented edge $e\in\mcE^{(N)}$,
\[
\mathcal A_e^{(N)}
=
\xi_e^{(N),\mu}\,A_\mu\bigl(\Xi^{(N)}(s(e))\bigr)
+O(a_N^2)
\]
uniformly in $e$.
Consequently,
\[
U_3^{(N)}(e)
=
I+i g_3\,\xi_e^{(N),\mu}A_\mu\bigl(\Xi^{(N)}(s(e))\bigr)
+O(a_N^2)
\]
uniformly in $e$.
\end{theorem}
\begin{proof}
Fix $e=(x\to y)\in\mcE^{(N)}$ and write $z_x:=\Xi^{(N)}(x)$ and
$\xi_e:=\Xi^{(N)}(y)-\Xi^{(N)}(x)$.
By the $C^2$ regularity of $\mathbf c$ and the color-field approximation,
\[
\mathbf c_y^{(N)}-\mathbf c_x^{(N)}
=
\xi_e^\mu\,\partial_\mu\mathbf c(z_x)+O(a_N^2)
\]
uniformly in $e$.
Substituting this into Definition~\ref{def:transport} gives
\begin{align*}
\mathcal A_e^{(N)}
&=
\frac12\sum_{a=1}^8
\Re\!\bigl((\mathbf c_x^{(N)})^\dagger\lambda_a(\mathbf c_y^{(N)}-\mathbf c_x^{(N)})\bigr)\lambda_a\\
&=
\frac12\sum_{a=1}^8
\Re\!\bigl(\mathbf c(z_x)^\dagger\lambda_a\,\xi_e^\mu\partial_\mu\mathbf c(z_x)\bigr)\lambda_a
+O(a_N^2)\\
&=
\xi_e^\mu A_\mu(z_x)+O(a_N^2).
\end{align*}
This proves the first claim.
Because $\|\Theta_e^{(N)}\|=g_3\|\mathcal A_e^{(N)}\|=O(a_N)$ uniformly,
the matrix exponential satisfies
\[
U_3^{(N)}(e)=I+i\Theta_e^{(N)}+O(a_N^2),
\]
which yields the second display.
\end{proof}

\begin{theorem}[Small-loop holonomy expansion]
\label{thm:small-loop-expansion}
Assume the admissible refinement regime of Definition~\ref{def:refinement}.
Let
\[
F_{\mu\nu}
:=
\partial_\mu A_\nu-\partial_\nu A_\mu-i g_3[A_\mu,A_\nu]
\]
be the curvature of the connection from
Theorem~\ref{thm:short-link-connection}.
Then for every plaquette $P\in\mathcal C_\square^{(N)}$,
\[
U_3^{(N)}(P)
=
\exp\!\Bigl(i g_3\,\Sigma_P^{(N),\mu\nu}F_{\mu\nu}(z_P^{(N)})
+O(a_N^3)\Bigr),
\]
with the remainder uniform in $P$.
In particular, for all sufficiently large $N$, $U_3^{(N)}(P)$ lies in the
principal logarithm domain and
\[
F_P^{(N)}
=
\frac{\Sigma_P^{(N),\mu\nu}}{A_P^{(N)}}\,F_{\mu\nu}(z_P^{(N)})+O(a_N).
\]
\end{theorem}
\begin{proof}
Write the oriented boundary of $P$ as $(e_1,e_2,e_3,e_4)$.
By Theorem~\ref{thm:short-link-connection},
\[
U_3^{(N)}(e_k)
=
\exp\!\Bigl(i g_3\,\xi_{e_k}^{(N),\mu}A_\mu(z_P^{(N)})+O(a_N^2)\Bigr)
\]
uniformly in $k$ and in $P$.
The plaquette boundary is closed, so
\[
\sum_{k=1}^4 \xi_{e_k}^{(N)}=0.
\]
Applying the Baker--Campbell--Hausdorff formula twice and using the
shape-regularity bounds from Definition~\ref{def:refinement} yields
\[
\log U_3^{(N)}(P)
=
i g_3\,\Sigma_P^{(N),\mu\nu}
\bigl(\partial_\mu A_\nu-\partial_\nu A_\mu-i g_3[A_\mu,A_\nu]\bigr)(z_P^{(N)})
+O(a_N^3).
\]
This is the first claim.
Because $\|\Sigma_P^{(N)}\|=O(a_N^2)$, the leading term in the logarithm tends to
zero uniformly, so for all sufficiently large $N$ the spectrum of $U_3^{(N)}(P)$
stays inside the principal logarithm domain.
Dividing the logarithm by $i g_3 A_P^{(N)}$ gives the final display.
\end{proof}

\begin{corollary}[Wilson plaquette density expansion]
\label{cor:wilson-density}
Under the hypotheses of Theorem~\ref{thm:small-loop-expansion}, each plaquette
term
\[
s_P^{(N)}:=1-\frac13\Re\tr U_3^{(N)}(P)
\]
satisfies
\[
s_P^{(N)}
=
\frac{g_3^2}{6}\,
\tr\!\Bigl(\bigl(\Sigma_P^{(N),\mu\nu}F_{\mu\nu}(z_P^{(N)})\bigr)^2\Bigr)
+r_P^{(N)},
\]
with a uniform remainder bound
\[
|r_P^{(N)}|\le C a_N^5.
\]
\end{corollary}
\begin{proof}
By Theorem~\ref{thm:small-loop-expansion},
\[
U_3^{(N)}(P)=\exp\!\bigl(i g_3 B_P^{(N)}+O(a_N^3)\bigr),
\qquad
B_P^{(N)}:=\Sigma_P^{(N),\mu\nu}F_{\mu\nu}(z_P^{(N)}).
\]
Since $\|B_P^{(N)}\|=O(a_N^2)$ uniformly and $\tr(B_P^{(N)})=0$,
Taylor expansion of the exponential gives
\[
\tr U_3^{(N)}(P)
=
3-\frac{g_3^2}{2}\tr\!\bigl((B_P^{(N)})^2\bigr)+O(a_N^5),
\]
where the linear term vanishes because $B_P^{(N)}$ is traceless Hermitian and the
order-$a_N^3$ correction in the exponent contributes only at order $a_N^5$ to
the real part of the trace.
Rearranging proves the claim.
\end{proof}

\begin{theorem}[Continuum limit of the $SU(3)$ Wilson action]
\label{thm:continuum-limit-ym}
Assume the admissible refinement regime of Definition~\ref{def:refinement}.
Then the discrete Wilson action from Definition~\ref{def:action} satisfies
\[
S_{\mathrm W}^{(N)}
\longrightarrow
\frac14\int_M \tr\!\bigl(F_{\mu\nu}F^{\mu\nu}\bigr)\,d\mathrm{vol}_g
\qquad\text{as }N\to\infty.
\]
\end{theorem}
\begin{proof}
By Corollary~\ref{cor:wilson-density},
\[
S_{\mathrm W}^{(N)}
=
\frac{6}{g_3^2}\sum_{P\in\mathcal C_\square^{(N)}} s_P^{(N)}
=
\sum_{P\in\mathcal C_\square^{(N)}}
\tr\!\Bigl(\bigl(\Sigma_P^{(N),\mu\nu}F_{\mu\nu}(z_P^{(N)})\bigr)^2\Bigr)
+\frac{6}{g_3^2}\sum_{P\in\mathcal C_\square^{(N)}} r_P^{(N)}.
\]
The remainder sum vanishes: by Definition~\ref{def:refinement},
$A_P^{(N)}\asymp a_N^2$, so $|\mathcal C_\square^{(N)}|=O(a_N^{-4})$, and hence
\[
\sum_{P\in\mathcal C_\square^{(N)}} |r_P^{(N)}|
\le C a_N^5 |\mathcal C_\square^{(N)}|
=O(a_N)\to0.
\]
For the leading term, write
\[
\tr\!\Bigl(\bigl(\Sigma_P^{(N),\mu\nu}F_{\mu\nu}(z_P^{(N)})\bigr)^2\Bigr)
=
\Sigma_P^{(N),\mu\nu}\Sigma_P^{(N),\rho\sigma}\,
\tr\!\bigl(F_{\mu\nu}(z_P^{(N)})F_{\rho\sigma}(z_P^{(N)})\bigr).
\]
Apply the bivector density convergence from Definition~\ref{def:refinement} with
\[
f_{\mu\nu\rho\sigma}(z)
:=
\tr\!\bigl(F_{\mu\nu}(z)F_{\rho\sigma}(z)\bigr).
\]
Since $F_{\mu\nu}$ is antisymmetric in $(\mu,\nu)$,
\begin{align*}
\sum_{P\in\mathcal C_\square^{(N)}}
\tr\!\Bigl(\bigl(\Sigma_P^{(N),\mu\nu}F_{\mu\nu}(z_P^{(N)})\bigr)^2\Bigr)
&\longrightarrow
\frac18\int_M
\bigl(g^{\mu\rho}g^{\nu\sigma}-g^{\mu\sigma}g^{\nu\rho}\bigr)
\tr(F_{\mu\nu}F_{\rho\sigma})\,d\mathrm{vol}_g\\
&=
\frac14\int_M \tr(F_{\mu\nu}F^{\mu\nu})\,d\mathrm{vol}_g.
\end{align*}
Combining the leading term with the vanishing remainder proves the theorem.
\end{proof}

\subsection{Quantization and partition function}

The data package $\mathcal Q$ provides a well-defined action
$S=S_{\mathrm W}+S_{\mathrm F}+S_\phi$
on the finite-dimensional field variables
$\{\psi_x,\bar\psi_x,\phi_x\}_{x\in\mcN}$.
The path integral over these fields,
\[
Z=\int \Bigl[\prod_{x\in\mcN}d\psi_x\,d\bar\psi_x\,d\phi_x\Bigr]
\exp\!\bigl(-S_{\mathrm W}-S_{\mathrm F}-S_\phi\bigr),
\]
is a finite-dimensional ordinary integral over finite-dimensional fermionic
Grassmann variables and the scalar sector $\R^{|\mcN|}$, evaluated in the
background of the explicit finite link configuration $\{U_3(e)\}_{e\in\mcE}$.
If one wishes to promote the link variables themselves to dynamical integration
variables, gauge-fixing via Faddeev--Popov is straightforward in the finite
setting.
The resulting $Z$ is well-defined without any ultraviolet regularization;
all divergences that would appear in a continuum QFT are absent because the
lattice is finite.
Evaluating $Z$ numerically for realistic run data and comparing with continuum
expectations is a natural next step.

\section{Conclusions}
\label{sec:conclusions}

We have constructed a complete, finite, and $SU(3)$ gauge-invariant lattice QFT
data package $\mathcal Q$ from the discrete trajectory data of a viscous
multi-particle system.
The construction proceeds in six conceptual steps:
\begin{enumerate}[leftmargin=1.5em]
\item Build the spacetime lattice $(\mcN,\mcE)$ from the recorded sites,
      worldline links, and same-time interaction links.
\item Compute the row-normalized viscous force $\mathbf F^{\mathrm{visc}}(x)$
      at each site from the interaction-link weights.
\item Form complex color vectors $\mathbf c_x\in\C^3$ by modulating each force
      component with a de~Broglie kinematic phase and normalizing.
\item Define the $SU(3)$ link transport $U_3(e)=\exp(i\Theta_e)$ by projecting
      the short-link color increment
      $\mathbf c_x^\dagger\lambda_a(\mathbf c_y-\mathbf c_x)$ onto the Lie
      algebra $\mathfrak{su}(3)$ via the Gell-Mann basis.
\item Build the canonical plaquette family $\mathcal C_\square$, the area
      bivectors $\Sigma_P$, the discrete plaquette holonomies $U_3(P)$, and
      the Wilson action $S_{\mathrm W}$.
\item Prove the thermodynamic mean-field limit for both link and plaquette
      observables, and under an admissible refinement regime recover the
      continuum Yang--Mills action.
\end{enumerate}
All objects are explicit and finite (Theorem~\ref{thm:package}).
The gauge group $SU(3)$ is not imposed but emerges from the three-dimensional
viscous dynamics together with the Gell-Mann projection.
In the stationary thermodynamic regime, the same interaction graph has a
deterministic mean-field limit in the sense of
Theorems~\ref{thm:edge-measure-limit} and~\ref{thm:plaquette-measure-limit}, so
scaled link and plaquette sums become mean-field interaction integrals without
changing the graph rule.
In the geometric small-loop regime, Theorems~\ref{thm:short-link-connection},
\ref{thm:small-loop-expansion}, and~\ref{thm:continuum-limit-ym} identify the
continuum connection and prove that the discrete Wilson action converges to
\[
\frac14\int_M \tr(F_{\mu\nu}F^{\mu\nu})\,d\mathrm{vol}_g.
\]

Two major extensions are deferred to companion papers:
(i) path-integral quantization and numerical evaluation of $Z$;
(ii) axiom verification --- checking that $\mathcal Q$ satisfies the
Osterwalder-Schrader axioms adapted to the finite lattice.

\bibliographystyle{plain}
\bibliography{references}

\end{document}
